\begin{trapbox}
\textbf{误区一：所有焦点弦都是通径。} 除了过焦点，还必须垂直于对应的轴。
\begin{center}\begin{tikzpicture}[font=\small,>=stealth]
\draw[black!30,rounded corners=2pt] (0,0) rectangle (14.4,8.2);
\node[anchor=west,font=\small\bfseries] at (.5,7.72) {同一个椭圆、同一个焦点：斜焦点弦不能套通径公式};
\begin{scope}[shift={(3.6,4.3)},scale=.43]
\draw[blue,very thick] (0,0) ellipse (5 and 3);\draw[black!35] (-5.5,0)--(5.5,0);\draw[orange!85!black,very thick] (4.0000000000,-1.8000000000)--(4.0000000000,1.8000000000);\node[fill=orange!85!black,circle,inner sep=1.4pt] at (4.0000000000,-1.8000000000) {};\node[fill=orange!85!black,circle,inner sep=1.4pt] at (4.0000000000,1.8000000000) {};\fill (4,0) circle (3pt) node[above left=6pt] {$F$};\end{scope}
\node at (3.6,1.3) {通径 $x=4$：$L=18/5$};
\begin{scope}[shift={(10.8,4.3)},scale=.43]
\draw[blue,very thick] (0,0) ellipse (5 and 3);\draw[black!35] (-5.5,0)--(5.5,0);\draw[orange!85!black,very thick] (1.0694232263,-2.9305767737)--(4.8129297149,0.8129297149);\node[fill=orange!85!black,circle,inner sep=1.4pt] at (1.0694232263,-2.9305767737) {};\node[fill=orange!85!black,circle,inner sep=1.4pt] at (4.8129297149,0.8129297149) {};\fill (4,0) circle (3pt) node[above left=6pt] {$F$};\end{scope}
\node at (10.8,1.3) {斜弦 $y=x-4$：$L=90/17$};

\node[anchor=west] at (.7,.55) {同尺度；均为 $x^2/25+y^2/9=1$；只有左图弦垂直于长轴};
\end{tikzpicture}\end{center}

右图仍过 $F=(4,0)$，但斜率为 $1$，整弦长 $90/17\ne18/5$。
回验：代入 $y=x-4$ 得 $34x^2-200x+175=0$，两横坐标差为 $45\sqrt2/17$，
乘 $\sqrt{1+1^2}$ 得 $90/17$；不能套通径式。

\medskip
\textbf{误区二：只取正端点坐标，漏乘2。}
横向模型上端点的纵坐标是半通径；整弦应算 $y_+-y_-$，不能只写 $b^2/a$ 或 $p$。

\medskip
\textbf{误区三：把 $a^2$ 当 $a$，或把 $p$ 当成焦点横坐标。}
有心曲线公式分母为半轴长 $a$；抛物线 $y^2=8x$ 中 $2p=8$，故 $p=4$、焦点横坐标 $2$、通径 $8$。
双曲线的 $a$ 要读正项分母的正平方根，不能按分母大小判断。

\medskip
\textbf{误区四：不在原点或改变开口就不能用长度公式。}
长度关系在平移、旋转、反射后保持；改变的是坐标表达。
应辨认轴与焦点，再确定通径方向，而不是强行固定为竖直线。


\medskip
\textbf{误区五：把“焦点轴”和“通径所在直线”弄成同一条。}
横向曲线的焦点轴是 $y=0$，但通径是 $x=\pm c$（有心曲线）
或 $x=p/2$（右开口抛物线）、$x=-p/2$（左开口抛物线）；
纵向模型要把 $x,y$ 的职责互换。
一眼判定：\textbf{通径一定与焦点轴垂直}。

\medskip
\textbf{误区六：双曲线里仍套椭圆的 $c^2=a^2-b^2$。}
椭圆是 $c^2=a^2-b^2$，双曲线是 $c^2=a^2+b^2$。
只要把对应的焦点坐标代回曲线，最终两者都得
$y^2=b^4/a^2$，但中途不可以混用 $c$ 的关系。

\medskip
\textbf{误区七：把“半轴长”认成方程分母的平方数，或误读 $p$。}
若 $x^2/25+y^2/9=1$，则 $a=5,b=3$，\textbf{不是} $a=25,b=9$。
若 $y^2=8x$，则 $2p=8$、焦点横坐标是 $p/2=2$；
三者 $2p,p,p/2$ 分别表示\textbf{整通径、焦准距、焦点到顶点的距离}。

\medskip
\textbf{误区八：把椭圆条件直接搬给双曲线。}
横向椭圆须 $a>b>0$；横向双曲线仅要求 $a>0,b>0$，
不存在“一定 $a>b$”的限制。双曲线中 $a$ 是实半轴长，
对应的是标准方程\textbf{正项}分母的平方根。

\end{trapbox}
